\documentclass[a4paper,11pt]{article}

\usepackage[a4paper,margin=2cm]{geometry}
\usepackage[T1]{fontenc}
\usepackage[utf8]{inputenc}
\usepackage[ngerman,english]{babel}
\usepackage{lmodern}
\usepackage{microtype}
\usepackage[table]{xcolor}
\usepackage{parskip}
\usepackage{enumitem}
\usepackage{array}
\usepackage{longtable}
\usepackage[most]{tcolorbox}
\usepackage{titlesec}
\usepackage[hidelinks]{hyperref}

\definecolor{HeaderBlueDark}{RGB}{70,110,170}
\definecolor{RowBlueLight}{RGB}{235,243,255}
\definecolor{RowBlueDark}{RGB}{220,233,250}
\definecolor{SoftGray}{RGB}{90,90,90}

\setlength{\parindent}{0pt}
\setlist[itemize]{leftmargin=1.4em,itemsep=0.35em,topsep=0.35em}
\renewcommand{\arraystretch}{1.35}
\setlength{\tabcolsep}{6pt}
\linespread{1.06}

\titleformat{\section}[block]
{\Large\bfseries\color{HeaderBlueDark}}
{}{0pt}{}
\titlespacing{\section}{0pt}{1.3em}{0.8em}

\titleformat{\subsection}[block]
{\large\bfseries\color{HeaderBlueDark}}
{}{0pt}{}
\titlespacing{\subsection}{0pt}{1.0em}{0.6em}

\newtcolorbox{exbox}{enhanced,breakable,colback=RowBlueLight,colframe=RowBlueDark,boxrule=0.4pt,arc=1.5mm,left=1.2mm,right=1.2mm,top=1mm,bottom=1mm,title={Beispiele \textnormal{(Examples)}},fonttitle=\bfseries,colbacktitle=RowBlueDark,coltitle=black}

\NewDocumentEnvironment{grammarentry}{mm+b}{%
    \begin{tcolorbox}[enhanced,breakable,colback=white,colframe=HeaderBlueDark,boxrule=0.6pt,arc=1.5mm,left=1.5mm,right=1.5mm,top=1mm,bottom=1mm,colbacktitle=HeaderBlueDark,coltitle=white,fonttitle=\bfseries,title={#1\hfill\normalfont\footnotesize #2}]
    #3
    \end{tcolorbox}
}{}

\newcommand{\GermanText}[1]{\textbf{Deutsch: }#1\par}
\newcommand{\EnglishText}[1]{{\small\color{SoftGray}\textbf{English: }#1\par}}
\newcommand{\ExampleItem}[2]{\item \textbf{#1}\\{\small\color{SoftGray}#2}}

\title{\glqq wie\grqq{} als subordinierende Konjunktion}
\date{}

\begin{document}
\selectlanguage{ngerman}
\maketitle
\tableofcontents
\newpage

\section{Grundfunktion}

\begin{grammarentry}{wie}{subordinierende Konjunktion / Subjunktion}
\GermanText{Als subordinierende Konjunktion leitet \textbf{wie} einen Nebensatz ein. In dieser Funktion bedeutet es meistens \glqq so, wie\grqq{} oder englisch \glqq as\grqq.}
\EnglishText{As a subordinating conjunction, \textbf{wie} introduces a subordinate clause. In this function, it usually means ``the way that'' or ``as.''}

\GermanText{Der mit \textbf{wie} eingeleitete Nebensatz beschreibt häufig die Art und Weise, in der etwas geschieht oder gilt.}
\EnglishText{The subordinate clause introduced by \textbf{wie} often describes the manner in which something happens or is valid.}
\end{grammarentry}

\section{Wortstellung}

\begin{grammarentry}{Verb am Ende}{Nebensatzregel}
\GermanText{Im \textbf{wie}-Nebensatz steht das konjugierte Verb normalerweise am Ende.}
\EnglishText{In a subordinate clause introduced by \textbf{wie}, the conjugated verb normally stands at the end.}

\GermanText{Steht der Nebensatz vor dem Hauptsatz, besetzt er die erste Position des Gesamtsatzes. Deshalb folgt im Hauptsatz sofort das konjugierte Verb.}
\EnglishText{When the subordinate clause comes before the main clause, it occupies the first position of the whole sentence. Therefore, the conjugated verb follows immediately in the main clause.}
\end{grammarentry}

\begin{exbox}
\begin{itemize}
\ExampleItem{Wie es bisher erklärt worden ist, brauchen wir die richtigen Nachrichten.}{As it has been explained so far, we need the correct information.}
\ExampleItem{Wie du weißt, wohne ich in Wien.}{As you know, I live in Vienna.}
\ExampleItem{Wir machen es so, wie der Lehrer es erklärt hat.}{We do it the way the teacher explained it.}
\end{itemize}
\end{exbox}

\section{Analyse des Beispielsatzes}

\begin{grammarentry}{Wie es bisher erklärt worden ist}{modaler Nebensatz}
\GermanText{In dem Satz \glqq Wie es bisher erklärt worden ist, brauchen wir die richtigen Nachrichten, um eine praktische Demokratie zu haben.\grqq{} leitet \textbf{wie} den Nebensatz \glqq wie es bisher erklärt worden ist\grqq{} ein.}
\EnglishText{In the sentence ``As it has been explained so far, we need the correct information in order to have a practical democracy,'' \textbf{wie} introduces the subordinate clause ``wie es bisher erklärt worden ist.''}

\GermanText{Der Nebensatz gibt an, entsprechend welcher Erklärung die folgende Aussage gemacht wird. Er hat deshalb eine modale Funktion.}
\EnglishText{The subordinate clause indicates according to which explanation the following statement is made. It therefore has a modal function.}
\end{grammarentry}

\rowcolors{2}{RowBlueLight}{RowBlueDark}
\begin{longtable}{>{\bfseries}p{4.1cm}|p{5.1cm}|p{5.1cm}}
\rowcolor{HeaderBlueDark}
\color{white}\textbf{Teil} &
\color{white}\textbf{Grammatische Funktion} &
\color{white}\textbf{English} \\
wie & subordinierende Konjunktion; leitet den Nebensatz ein & subordinating conjunction; introduces the subordinate clause \\
es & Subjekt des Nebensatzes & subject of the subordinate clause \\
bisher & Temporaladverb & temporal adverb: so far \\
erklärt worden ist & Prädikat im Perfekt Passiv; \textbf{ist} steht am Ende & predicate in the present perfect passive; \textbf{ist} is at the end \\
\end{longtable}

\section{Bedeutung: Art und Weise}

\begin{grammarentry}{modaler Gebrauch}{so, wie / as}
\GermanText{\textbf{wie} verbindet den Inhalt des Nebensatzes mit der Art und Weise des Hauptsatzes. Der Nebensatz beantwortet sinngemäß die Frage \glqq Wie?\grqq{} oder \glqq Auf welche Weise?\grqq.}
\EnglishText{\textbf{wie} connects the content of the subordinate clause with the manner of the main clause. In meaning, the subordinate clause answers the question ``How?'' or ``In what way?''}
\end{grammarentry}

\begin{exbox}
\begin{itemize}
\ExampleItem{Mach es, wie ich es dir gezeigt habe.}{Do it the way I showed you.}
\ExampleItem{Alles geschah, wie wir es erwartet hatten.}{Everything happened as we had expected.}
\ExampleItem{Wie bereits erwähnt wurde, beginnt der Kurs am Montag.}{As was already mentioned, the course begins on Monday.}
\end{itemize}
\end{exbox}

\section{Korrelat \glqq so\grqq}

\begin{grammarentry}{so ..., wie ...}{häufige Struktur}
\GermanText{Im Hauptsatz kann das Korrelat \textbf{so} stehen. Es verweist auf den folgenden \textbf{wie}-Nebensatz.}
\EnglishText{The correlating word \textbf{so} may appear in the main clause. It points forward to the subordinate clause introduced by \textbf{wie}.}

\GermanText{Struktur: \textbf{etwas so machen, wie ...}}
\EnglishText{Structure: \textbf{to do something the way that ...}}
\end{grammarentry}

\begin{exbox}
\begin{itemize}
\ExampleItem{Ich mache es so, wie du es vorgeschlagen hast.}{I do it the way you suggested.}
\ExampleItem{Der Text ist so geschrieben, wie es die Aufgabe verlangt.}{The text is written in the way the task requires.}
\end{itemize}
\end{exbox}

\section{Abgrenzung zum Fragewort}

\begin{grammarentry}{kein Fragewort}{in diesem Gebrauch}
\GermanText{In \glqq Wie es bisher erklärt worden ist\grqq{} ist \textbf{wie} kein Frageadverb. Es stellt keine Frage, sondern verbindet den Nebensatz mit dem Hauptsatz.}
\EnglishText{In ``Wie es bisher erklärt worden ist,'' \textbf{wie} is not an interrogative adverb. It does not ask a question; it connects the subordinate clause with the main clause.}
\end{grammarentry}

\rowcolors{2}{RowBlueLight}{RowBlueDark}
\begin{longtable}{>{\bfseries}p{3.2cm}|p{5.7cm}|p{5.4cm}}
\rowcolor{HeaderBlueDark}
\color{white}\textbf{Gebrauch} &
\color{white}\textbf{Beispiel} &
\color{white}\textbf{Erklärung} \\
Frageadverb & Wie funktioniert das? & direkte Frage; \textbf{wie} fragt nach der Art und Weise \\
Subordinierende Konjunktion & Wie du weißt, ist das wichtig. & \textbf{wie} leitet einen Nebensatz mit Verbendstellung ein \\
Vergleichswort & Er ist so groß wie sein Bruder. & Vergleich ohne vollständigen Nebensatz \\
\end{longtable}

\section{Kommasetzung}

\begin{grammarentry}{Komma}{Trennung von Haupt- und Nebensatz}
\GermanText{Ein vollständiger \textbf{wie}-Nebensatz wird durch ein Komma vom Hauptsatz getrennt.}
\EnglishText{A complete subordinate clause introduced by \textbf{wie} is separated from the main clause by a comma.}

\GermanText{Ein bloßer Vergleich ohne vollständigen Nebensatz braucht dagegen normalerweise kein Komma: \glqq Er ist so groß wie sein Bruder.\grqq}
\EnglishText{A simple comparison without a complete subordinate clause normally does not require a comma: ``He is as tall as his brother.''}
\end{grammarentry}

\section{Wichtiger Bedeutungsunterschied}

\begin{grammarentry}{wie und als}{Vergleich}
\GermanText{Bei Gleichheit verwendet man normalerweise \textbf{wie}: \glqq so groß wie\grqq. Bei Ungleichheit verwendet man \textbf{als}: \glqq größer als\grqq. Diese Vergleichsfunktion ist von dem hier behandelten subordinierenden Gebrauch zu unterscheiden.}
\EnglishText{For equality, German normally uses \textbf{wie}: ``as tall as.'' For inequality, it uses \textbf{als}: ``taller than.'' This comparative function must be distinguished from the subordinating use discussed here.}
\end{grammarentry}

\section{Zusammenfassung}

\begin{grammarentry}{Merksatz}{kurz und wichtig}
\GermanText{Als subordinierende Konjunktion bedeutet \textbf{wie} meistens \glqq so, wie\grqq{} oder \glqq as\grqq. Es leitet einen Nebensatz ein, und das konjugierte Verb steht am Ende.}
\EnglishText{As a subordinating conjunction, \textbf{wie} usually means ``the way that'' or ``as.'' It introduces a subordinate clause, and the conjugated verb stands at the end.}

\GermanText{Beispiel: \textbf{Wie du weißt, ist Deutsch nicht immer leicht.}}
\EnglishText{Example: \textbf{As you know, German is not always easy.}}
\end{grammarentry}

\end{document}
